\chapter{Game Version, Names, and Guide Conventions}
\label{written-conventions}

\emph{Dragon Quest V} begins aboard a ship, with a young boy travelling
alongside his father, Papas. The first visits to Santa Rosa and Alcapa
establish a comfortable rhythm of talking to townspeople, searching houses,
buying equipment, and venturing into the countryside. As the journey grows,
the party gains new companions and new ways to travel. Familiar places
remain part of the adventure, and return visits give them fresh significance.

The game rewards attention to small details. A drawer can hold a useful
healing item. A conversation can point toward the next destination. A new
key can make several earlier stops worth revisiting. Keeping track of these
connections makes exploration satisfying and helps the party arrive ready
for the next challenge.

\section{Game Version}

\emph{Dragon Quest V: Tenkuu no Hanayome} is Enix's 1992 role-playing game
for the Super Famicom / SNES.

The subtitle is commonly rendered as \emph{The Bride of the Sky} or
\emph{Bride of Heaven}. The names \emph{Dragon Quest V}, \emph{Dragon
Quest 5}, and \emph{DQ5} identify the same game.

Exploration, party development, and monster recruitment are central to the
adventure. Several features shape the journey:
\begin{itemize}
  \item The normal active battle party has three members. Choosing their
    roles and equipment is an important part of preparing for a dungeon.
  \item The wagon carries additional companions. Access to those companions
    depends on the location and the party's current circumstances.
  \item Each recruitable monster species has its own growth, abilities,
    equipment options, and maximum level.
  \item Keys and vehicles expand the accessible world throughout the game.
    A new travel ability can make an earlier destination worth another visit.
\end{itemize}

\section{Names}

\subsection*{One place, several names}

Alcapa, Alkapa, and Arukapa identify the same town. Lenoire and Linoire
are alternate spellings of the same castle. Common place names include
the following:

\begin{guidetable}
\begin{center}
\begin{tabular}{@{}p{0.28\textwidth}p{0.64\textwidth}@{}}
\toprule
\GuideTableHeading{Primary name} & \GuideTableHeading{Alternate names} \\
\midrule
Santa Rosa & Santa Rose, Santorez \\
Alcapa & Alkapa, Arukapa \\
Lenoire Castle & Linoire Castle, Lenule, Linule \\
Oracleberry & Oracle Valley \\
Salabona & Sarabona \\
Granvania & Granbania, Glanvania \\
El Heven & Elheaven \\
Torokko Mines & Trolley Cave \\
\bottomrule
\end{tabular}
\end{center}
\end{guidetable}

Surroundings provide another useful way to identify a place. Its position
relative to a port, mountain pass, river, or nearby town helps confirm the
match. Dungeon landmarks, floor connections, and the current objective
provide similar clues when a place has an unfamiliar name.

\subsection*{People and companion names}

The \emph{Hero} is the player-named protagonist who travels with Papas at
the beginning of the adventure.

Monster companions have both a species and an individual name. A Slime
Knight's species determines its growth, spells, and equipment options.
Individual companions can have their own nicknames, and the roster may
contain more than one member of a species.

A character's appearance, location, and abilities help establish its
identity. These details are useful when several enemy or companion names
resemble one another.

\subsection*{Spells and equipment}

\emph{Blaze}, \emph{Firebal}, and \emph{Bang} belong to separate spell
families. MP cost, target pattern, and effect help identify a spell and
its role in battle.

Weapons and armor have individual names, power values, eligible users,
and special properties. Certain rings and staves also have item-use
effects alongside their equipment properties.

